% Dowody według źródła 2, ze sprawdzeniem w źródle 1.
\poczatekwykladu{7}

\subsection{Jednoznaczność rozkładu i związek z gęstością}

\begin{twierdzenie}[Dystrybuanta wyznacza rozkład]
\label{tw:dystrybuanta-wyznacza-rozklad}
Jeżeli rozkłady $\mu$ i $\nu$ na $\R$ mają tę samą dystrybuantę,
to $\mu(A)=\nu(A)$ dla każdego $A\in\B(\R)$.
\end{twierdzenie}

\begin{proof}
Załóżmy, że $F_\mu(t)=F_\nu(t)$ dla każdego $t\in\R$. Zdefiniujmy
\[
  \mathcal A=\{A\in\B(\R):\mu(A)=\nu(A)\},\qquad
  \mathcal P=\{(-\infty,t]:t\in\R\}.
\]
Rodzina $\mathcal P$ jest $\pi$-układem, ponieważ przecięcie dwóch
takich półprostych ma prawy koniec w mniejszym z ich prawych końców.
Generuje ona $\B(\R)$, gdyż z półprostych można otrzymać przedziały
$(a,b]$ przez różnicę, a stąd przedziały otwarte przez przeliczalne sumy.
Ponadto $\mathcal P\subseteq\mathcal A$, gdyż dla każdego $t\in\R$
\[
  \mu((-\infty,t])=F_\mu(t)=F_\nu(t)=\nu((-\infty,t]).
\]

Sprawdzamy, że $\mathcal A$ jest $\lambda$-układem.
\begin{enumerate}
  \item Mamy $\R\in\mathcal A$, ponieważ $\mu(\R)=\nu(\R)=1$.
  \item Jeżeli $A,B\in\mathcal A$ i $A\subseteq B$, to
    \[
      \mu(B\setminus A)=\mu(B)-\mu(A)=\nu(B)-\nu(A)=\nu(B\setminus A),
    \]
    zatem $B\setminus A\in\mathcal A$.
  \item Jeżeli $A_n\in\mathcal A$ oraz $A_n\subseteq A_{n+1}$
    dla każdego $n$, to z ciągłości miar od dołu wynika
    \[
      \mu\!\left(\bigcup_{n=1}^\infty A_n\right)
      =\lim_{n\to\infty}\mu(A_n)
      =\lim_{n\to\infty}\nu(A_n)
      =\nu\!\left(\bigcup_{n=1}^\infty A_n\right).
    \]
    Stąd $\bigcup_{n=1}^\infty A_n\in\mathcal A$.
\end{enumerate}
Twierdzenie o układach $\pi$ i $\lambda$ (twierdzenie~\ref{tw:pi-lambda})
daje $\B(\R)=\sigma(\mathcal P)\subseteq\mathcal A$,
więc $\mathcal A=\B(\R)$ i $\mu=\nu$.
\end{proof}

Na przedziałach półotwartych wartość rozkładu można odczytać bezpośrednio
z dystrybuanty. Dla $a<b$ zachodzi
\[
  \mu((a,b])=\mu((-\infty,b])-\mu((-\infty,a])=F_\mu(b)-F_\mu(a).
\]

\begin{fakt}[Dystrybuanta rozkładu absolutnie ciągłego]
\label{fakt:dystrybuanta-z-gestosci}
Rozkład zmiennej losowej $X$ jest absolutnie ciągły wtedy i tylko wtedy,
gdy istnieje nieujemna funkcja borelowska $g\colon\R\to[0,\infty)$
taka, że
\begin{equation}
  F_X(t)=\int_{-\infty}^t g(x)\,dx
  \qquad\text{dla każdego }t\in\R.
  \label{eq:dystrybuanta-calka-gestosci}
\end{equation}
Wówczas $g$ jest gęstością tego rozkładu.
\end{fakt}

\begin{proof}
Jeżeli $\mu_X$ jest absolutnie ciągły i ma gęstość $g$, to
\[
  F_X(t)=\mu_X((-\infty,t])=\int_{-\infty}^{t}g(x)\,dx
\]
na mocy definicji~\ref{def:rozklad-absolutnie-ciagly}.

Odwrotnie, załóżmy, że $g$ spełnia~\eqref{eq:dystrybuanta-calka-gestosci}.
Z nieujemności $g$ i ciągłości całki od dołu otrzymujemy
\[
  \int_{\R}g(x)\,dx
    =\lim_{n\to\infty}\int_{-\infty}^n g(x)\,dx
    =\lim_{n\to\infty}F_X(n)=1.
\]
Wobec tego $\nu(A)=\int_A g(x)\,dx$, $A\in\B(\R)$,
definiuje rozkład absolutnie ciągły. Jego dystrybuanta spełnia
$F_\nu(t)=\int_{-\infty}^t g(x)\,dx=F_X(t)$ dla każdego $t\in\R$.
Z twierdzenia~\ref{tw:dystrybuanta-wyznacza-rozklad} wynika $\nu=\mu_X$,
a więc $\mu_X$ ma gęstość $g$.
\end{proof}
